\section{Introduction}
\Acp{ANN} are the core of many modern applications, for example computer vision or natural language processing.
They try to mimic the neural behavior inside the brain.
Traditionally, a \ac{ANN} is based on the multi layer perceptron model, where neurons are grouped in layers.
Each neuron computes a real-valued numeric value based on its weighted input values.
% TODO search for source
The value is transformed by an activation function and then passed as output to the next layer.

\Acp{SNN} offer an alternative paradigm to traditional \acp{ANN}.
In contrast to \acp{ANN}, \Acp{SNN} process discrete, time-dependent signals, referred to as \textit{spikes}.
A neuron in an \Acp{SNN} does not respond to a single input but instead integrates multiple spikes over time~\cite{pietrzak_overview_2023}.
Once the so-called \textit{membrane potential} exceeds a certain threshold, the neuron emits an outgoing spike.
This mechanisms mimics the behavior of the brain in a more biologically plausible way compared to the perceptron model~\cite{maass_networks_1997}.
As a result, \acp{SNN} become event-driven: information is only transmitted when spikes are emitted.
This leads to asynchronous computations, unlike synchronous computations in \acp{ANN}~\cite{pietrzak_overview_2023}.

Beside more plausibility from the biological point of view, \acp{SNN} are considered potentially more energy-efficient than \acp{ANN}~\cite{pietrzak_overview_2023}.
However, their practical adoption remains limited.
One major reason is that their computational model aligns best with neuromorphic hardware, which is designed specifically to handle time-dependent and sparse activity patterns.
Unfortunately, such hardware is not yet widely available~\cite{li_towards_2024}.

In contrast, \acp{GPU} are widespread and there are a several optimizations for \acp{ANN} running on \acp{GPU} (e.g., ~\cite{alevi_brian2cuda_2022, vieth_accelerating_2024}).
However, the asynchronous and event-driven nature of \acp{SNN} is fundamentally different from the dense, synchronous computation pattern \acp{GPU} are designed for.
This raises the question of whether \acp{GPU} can serve as an effective platform for training and executing \acp{SNN} — despite this mismatch.

This paper investigates this question by focusing on three main aspects:
\begin{enumerate}
    \item What are the architectural differences between \acp{SNN} and traditional \acp{ANN}?
    \item Which methodologies can be used to train SNNs on \acp{GPU}?
    \item How do \acp{SNN} perform on \acp{GPU} compared to \acp{ANN}?
\end{enumerate}

The paper is based on a literature review.
The main focus is set on papers which contain information about running \acp{SNN} on \acp{GPU} and corresponding frameworks like Brian or PymoNNtorch.

The paper is structured as follows:
% TODO rewrite structure section with refs
In the following section, the fundamentals of \acp{SNN} architecture will be presented and compared to \acp{ANN}, with a particular focus on neuron models, activation mechanisms, and signal propagation.
Next, the paper will discuss methodologies for training and running \acp{SNN} on \acp{GPU}, including their advantages and limitations.
After that, a comparison of \acp{SNN} and traditional \acp{ANN} will be conducted with a focus on performance aspects when executed on \acp{GPU}.
Finally, the paper will be summarized.

\newpage